\documentclass[12pt]{article}
\usepackage[spanish, es-lcroman]{babel}
\usepackage{amssymb}
\usepackage{enumitem}
\usepackage{fancyhdr}
\usepackage{amsmath}

\pagestyle{fancy}
\title{Taller No.2 de Cálculo}
\author{Carlos David Pablo Hernández}
\date{\today}

\setlength{\headheight}{20pt}

\fancyhead[L]{Carlos Pablo}
\fancyhead[C]{Taller No. 2}
\fancyhead[R]{\today}


\begin{document}

\maketitle
\section{} Para los siguientes conjuntos e intervalos, determine si poseen cota superior e inferior. Justifique su respuesta
indicando el supremo y el ínfimo (si existen) y proporcione al menos dos ejemplos de números que actúen como
cotas superiores y dos como cotas inferiores para cada caso.
\begin{itemize}
    \item $(-2, 4]$ \\
    El intervalo posee cotas superiores e inferiores. \\
    \textbf{Supremo:} 4 \\
    \textbf{Ínfimo:} -2\\
    \textbf{Ejemplos de Cotas Superiores:} 5, 6 \\
    \textbf{Ejemplos de Cotas Inferiores:} -3, -4 \\

     \item $(5, +\infty)$ \\
    El intervalo posee cotas inferiores. \\
    \textbf{Supremo:} No tiene supremo. \\
    \textbf{Ínfimo:} 5 \\
    \textbf{Ejemplos de Cotas Superiores:} No hay cotas superiores.\\
    \textbf{Ejemplos de Cotas Inferiores:} 4, 3 \\

     \item $B= \{-2, 7\} $ \\
    El conjunto posee cotas superiores e inferiores. \\
    \textbf{Supremo:} 7 \\
    \textbf{Ínfimo:} -2 \\
    \textbf{Ejemplos de Cotas Superiores:} 8, 9 \\
    \textbf{Ejemplos de Cotas Inferiores:} -3, -4 \\

     \item $[7, 0)$ \\
    El intervalo es un conjunto vacío, ya que no existen números reales
    que sean mayores o iguales a 7 y menores a 0 al mismo tiempo. 
    Por lo tanto, no tiene cotas superiores ni inferiores. \\

     \item $A=(2, 5)$ \\
    El intervalo posee cotas superiores e inferiores. \\
    \textbf{Supremo:} 5 \\
    \textbf{Ínfimo:} 2 \\
    \textbf{Ejemplos de Cotas Superiores:} 6, 7 \\
    \textbf{Ejemplos de Cotas Inferiores:} 1, 0 \\

    \item $A=(2, 5]$ \\
    El intervalo posee cotas superiores e inferiores. \\
    \textbf{Supremo:} 5 \\
    \textbf{Ínfimo:} 2 \\
    \textbf{Ejemplos de Cotas Superiores:} 6, 7 \\
    \textbf{Ejemplos de Cotas Inferiores:} 1, 0 \\

    \item $C=[2, 5]$ \\
    El intervalo posee cotas superiores e inferiores. \\
    \textbf{Supremo:} 5 \\
    \textbf{Ínfimo:} 2 \\
    \textbf{Ejemplos de Cotas Superiores:} 6, 7 \\
    \textbf{Ejemplos de Cotas Inferiores:} 1, 0 \\

    \item $D=[-3, +\infty)$ \\
    El intervalo posee cotas inferiores. \\
    \textbf{Supremo:} No tiene supremo. \\
    \textbf{Ínfimo:} -3 \\
    \textbf{Ejemplos de Cotas Superiores:} No tiene cotas superiores. \\
    \textbf{Ejemplos de Cotas Inferiores:} -4, -5 \\

    \item $E=[-1, 7) \cup   [9, 11] $ \\
    El intervalo posee cotas superiores e inferiores. \\
    \textbf{Supremo:} 11 \\
    \textbf{Ínfimo:} -1 \\
    \textbf{Ejemplos de Cotas Superiores:} 12, 13 \\
    \textbf{Ejemplos de Cotas Inferiores:} -2, -3 \\

    \item $F=\mathbb{R} $ \\
    El conjunto no posee cotas superiores ni inferiores. \\
    \textbf{Supremo:} No tiene supremo. \\
    \textbf{Ínfimo:} No tiene ínfimo. \\
    \textbf{Ejemplos de Cotas Superiores:} No tiene cotas superiores. \\
    \textbf{Ejemplos de Cotas Inferiores:} No tiene cotas inferiores. \\

    \item $G=\{-\infty \}$ \\
    El conjunto no se encuentra dentro de los números reales, por lo tanto, no tiene 
    cotas superiores ni inferiores. \\
    
    \item $H=(-\infty , 3] \cup   (4, +\infty ) $ \\
    El intervalo no posee cotas superiores ni inferiores. \\
    \textbf{Supremo:} No tiene supremo. \\
    \textbf{Ínfimo:} No tiene ínfimo. \\
    \textbf{Ejemplos de Cotas Superiores:} No tiene cotas superiores. \\
    \textbf{Ejemplos de Cotas Inferiores:} No tiene cotas inferiores. \\


\end{itemize}
\section{}
Responda y justifique su respuesta.

\begin{itemize}
    \item ¿Cuál de los siguientes conjuntos tiene a -7 como cota inferior? \\
    \begin{enumerate}[label=\alph*)]
        \item $(-\infty, 6) \cap [3, 7)$
        \item $(-\infty, -7)$
        \item $[2, 4) \cup (-8, 10)$
        \item $(-10, 10) \cap (-8, 0)$
    \end{enumerate}

    \textbf{Respuesta:} La opción a) es la única que tiene a $-7$ como cota inferior, 
    ya que la intersección entre los dos conjuntos es $[3, 6)$, y $-7$ es menor que cualquier 
    elemento de este conjunto. \\

    \item ¿Cuál de los siguientes conjuntos tiene a 5 como cota superior? \\
    \begin{enumerate}[label=\alph*)]
        \item $(5, +\infty)$
        \item $(-\infty, 5]$
        \item $[-5, 6)$
    \end{enumerate}

    \textbf{Respuesta:} La opción b) es la única que tiene a 5 como cota superior, 
    ya que todos los elementos del conjunto son menores o iguales a 5. \\

    \item ¿De cuál de los intervalos siguientes es cota inferior el número -1? \\
    \begin{enumerate}[label=\alph*)]
        \item $(-1, +\infty)$
        \item $(-\infty, -1)$
        \item $[-1, 1]$
    \end{enumerate}

    \textbf{Respuesta:} El $-1$ es cota inferior de los intervalos a) y c), 
    ya que en ambos casos, $-1$ es el límite inferior del intervalo, así que
    $-1$ es menor o igual a todos los elementos del intervalo. \\

\end{itemize}

\section{}
Halle la cota superior mínima y la cota inferior máxima (si existen) de los siguientes conjuntos. Decida también
cuáles de ellos poseen un elemento máximo y un elemento mínimo (es decir, en qué casos la cota superior mínima y
la cota inferior máxima pertenecen al conjunto).

\begin{enumerate}[label=\textit{\roman*})]
    
    \item $\{\frac{1}{n} : n \in \mathbb{N}\}$ \\
        \begin{align*}
            n= 1,\, \, \, \, \frac{1}{1} = 1 \\ \\
            n= 2,\, \, \, \, \frac{1}{2} = \frac{1}{2} \\ \\
            n= 3,\, \, \, \, \frac{1}{3} = \frac{1}{3} \\
            \vdots
        \end{align*}
    
    \textbf{Cota superior mínima:} 1  \\
    \textbf{Cota inferior máxima:} 0 \\
    \textbf{Elemento máximo:} 1 \\
    \textbf{Elemento mínimo:} No existe, ya que la cota inferior máxima
    no pertenece al conjunto.  \\

        \newpage

    \item $\{\frac{1}{n} : n \in \mathbb{Z},\ n \neq 0\}$ \\
        \begin{align*}
            n= 1,\, \, \, \, \frac{1}{1} = 1 \\ \\
            n= -1,\, \, \, \, \frac{1}{-1} = -1 \\ \\
            n= 2,\, \, \, \, \frac{1}{2} = \frac{1}{2} \\ \\
            n= -2,\, \, \, \, \frac{1}{-2} = -\frac{1}{2} \\ \\
            n= 3,\, \, \, \, \frac{1}{3} = \frac{1}{3} \\ \\
            n= -3,\, \, \, \, \frac{1}{-3} = -\frac{1}{3} \\ 
            \vdots
        \end{align*}

    \textbf{Cota superior mínima:} 1  \\
    \textbf{Cota inferior máxima:} -1 \\
    \textbf{Elemento máximo:} 1 \\
    \textbf{Elemento mínimo:} -1 \\


    \item $\{x : x = 0$ o $x = \frac{1}{n}$ para algún $n \in \mathbb{N}\}$ \\
        \begin{align*}
            n= 1,\, \, \, \, \frac{1}{1} = 1 \\ \\
            n= 2,\, \, \, \, \frac{1}{2} = \frac{1}{2} \\ 
            \vdots
        \end{align*}

    \textbf{Cota superior mínima:} 1  \\
    \textbf{Cota inferior máxima:} 0 \\
    \textbf{Elemento máximo:} 1 ($\frac{1}{1}$)\\
    \textbf{Elemento mínimo:} 0 (porque $x = 0$ está incluido en el conjunto)\\


    \item $\{x : 0 \leq x \leq \sqrt{2}$ y x es racional$\}$ \\ \\
    \textbf{Cota superior mínima:} $\sqrt{2}$  \\
    \textbf{Cota inferior máxima:} 0 \\
    \textbf{Elemento máximo:} No existe en el conjunto, porque $\sqrt{2}$ no es racional\\
    \textbf{Elemento mínimo:} 0 (porque $x = 0$ está incluido en el conjunto)\\

        \newpage

    \item $\{x : -x^2 + x + 1 \geq 0\}$ \\
        \begin{align*}
            -x^2 + x + 1 &\geq 0 \\ \\
            x^2 - x - 1 &\leq 0 \\ \\
            x^2 - x &\leq 1 \\ \\
            x^2 - x + \frac{1}{4} - \frac{1}{4} &\leq 1 \\ \\
            (x - \frac{1}{2})^2 - \frac{1}{4} &\leq 1 \\ \\
            (x - \frac{1}{2})^2 &\leq 1 + \frac{1}{4} \\ \\
            (x - \frac{1}{2})^2 &\leq \frac{5}{4} \\ \\
            \sqrt{(x - \frac{1}{2})^2} &\leq \sqrt{\frac{5}{4}} \\ \\
            |x - \frac{1}{2}| &\leq \frac{\sqrt{5}}{2} \\ \\
            -\frac{\sqrt{5}}{2} \leq \, &x - \frac{1}{2} \leq \frac{\sqrt{5}}{2} \\
            -\frac{\sqrt{5}}{2} + \frac{1}{2} \leq \, &x \leq \frac{\sqrt{5}}{2} + \frac{1}{2} \\ \\
            \frac{1 - \sqrt{5}}{2} \leq \,&x \leq \frac{1 + \sqrt{5}}{2}
        \end{align*}

    \textbf{Cota superior mínima:} $\frac{1 + \sqrt{5}}{2}$  \\
    \textbf{Cota inferior máxima:} $\frac{1 - \sqrt{5}}{2}$ \\
    \textbf{Elemento máximo:} $\frac{1 + \sqrt{5}}{2}$\\
    \textbf{Elemento mínimo:} $\frac{1 - \sqrt{5}}{2}$\\


    \item $\{x : x^2 + x - 1 < 0\}$ \\ 
        \begin{align*}
            x^2 + x - 1 &< 0 \\ \\
            x^2 + x &< 1 \\ \\
            (x + \frac{1}{2})^2 - \frac{1}{4} &< 1 \\ \\
            (x + \frac{1}{2})^2 &< 1 + \frac{1}{4} \\ \\
            (x + \frac{1}{2})^2 &< \frac{5}{4} \\ \\
            \sqrt{(x + \frac{1}{2})^2} &< \sqrt{\frac{5}{4}} \\ \\
            |x + \frac{1}{2}| &< \frac{\sqrt{5}}{2} \\ \\
            -\frac{\sqrt{5}}{2} < \, &x + \frac{1}{2} < \frac{\sqrt{5}}{2} \\
            -\frac{\sqrt{5}}{2} - \frac{1}{2} < \, &x < \frac{\sqrt{5}}{2} - \frac{1}{2} \\ \\
            \frac{-1 -\sqrt{5}}{2} <\, &x < \frac{-1 +\sqrt{5}}{2}
        \end{align*}

    \textbf{Cota superior mínima:} $\frac{-1 + \sqrt{5}}{2}$  \\
    \textbf{Cota inferior máxima:} $\frac{-1 - \sqrt{5}}{2}$ \\
    \textbf{Elemento máximo:} No existe en el conjunto, porque la cota superior
    mínima no pertenece al conjunto.\\
    \textbf{Elemento mínimo:} No existe en el conjunto, porque la cota inferior
    máxima no pertenece al conjunto.\\

    
    \item $\{x : x < 0$ y $x^2 + x - 1 < 0 \}$\\ \\
    Sabemos que el intervalo de 
        \[
            x^2 + x - 1 < 0
        \]
    es:    
        \[
            (\frac{-1 -\sqrt{5}}{2}, \frac{-1 +\sqrt{5}}{2})
        \]
    Entonces recortamos el intervalo hasta 0:
        \[
            (\frac{-1 -\sqrt{5}}{2}, 0)
        \]
    \textbf{Cota superior mínima:} $0$  \\
    \textbf{Cota inferior máxima:} $\frac{-1 - \sqrt{5}}{2}$ \\
    \textbf{Elemento máximo:} No existe en el conjunto, porque la cota superior
    mínima no pertenece al conjunto.\\
    \textbf{Elemento mínimo:} No existe en el conjunto, porque la cota 
    inferior máxima no pertenece al conjunto.\\

        \newpage

    \item $\{\frac{1}{n} + (-1)^n : n \in \mathbb{N}  \}$\\
        \begin{align*}
            n= 1,\, \, \, \, \frac{1}{1} + (-1)^1 = 0 \\ \\
            n= 2,\, \, \, \, \frac{1}{2} + (-1)^2 = \frac{3}{2} \\ \\
            n= 3,\, \, \, \, \frac{1}{3} + (-1)^3 = -\frac{2}{3} \\ \\
            n= 4,\, \, \, \, \frac{1}{4} + (-1)^4 = \frac{5}{4} \\ \\
            n= 5,\, \, \, \, \frac{1}{5} + (-1)^5 = -\frac{4}{5} \\ \\
            n= 6,\, \, \, \, \frac{1}{6} + (-1)^6 = \frac{7}{6} \\ \\
        \end{align*}    
     \textbf{Cota superior mínima:} $\frac{3}{2}$  \\
    \textbf{Cota inferior máxima:} $-1$ \\
    \textbf{Elemento máximo:} $\frac{3}{2}$\\
    \textbf{Elemento mínimo:} No existe en el conjunto, porque la 
    cota inferior máxima no pertenece al conjunto.\\

    \item \textbf{En el conjunto de los racionales cada racional $r$ tiene una única 
    representación irreducible de la forma $\frac{p}{q}$ donde $q$ es un 
    entero positivo. Luego a cada representación de este estilo se le borra
     el numerador y se cambia por 1, así obtenemos un nuevo conjunto de 
     números, determine el ínfimo y supremo de este conjunto.} \\
        El conjunto resultante es $\{\frac{1}{n} : n \in \mathbb{N}\}$, ya que cada racional $r$ se representa como $\frac{p}{q}$, y al cambiar el numerador por 1, obtenemos $\frac{1}{q}$, donde $q$ es un entero positivo. Por lo tanto, el conjunto resultante es el mismo que el conjunto del inciso i). \\
     \textbf{Supremo:} 1 \\
     \textbf{Ínfimo:} 0 \\
    \textbf{Elemento máximo:} $1$\\
    \textbf{Elemento mínimo:} No existe en el conjunto, porque el ínfimo
     no pertenece al conjunto.\\

\end{enumerate}

\section{}
Argumente, no es necesario demostrar matemáticamente. \\
\begin{itemize}
    \item \textbf{¿Por qué $\mathbb{N}$  no tiene cota superior pero sí inferior?} \\
    Porque dentro de los números naturales, no existe un número real que sea mayor
    a todos los elementos del conjunto; sin embargo, sí está acotado inferiormente, ya que 
    existen números dentro de los reales que son menores o iguales a todos los
    números naturales. \\

    \item \textbf{¿Por qué $\mathbb{Z}$  no está acotado superiormente?} \\
    Porque no existe un número real que sea mayor a todos los elementos del conjunto de los 
    enteros. \\
\end{itemize}
\end{document}